\chapter{Mass, propulsion, atmosphere, and reference frames}\label{ch:models}
The flight engine does not calculate every physical property itself. It asks separate model objects for mass, motor thrust, atmosphere, wind, gravity, and aerodynamic coefficients. \class{SimulationConditions} brings these services together for one scenario. \class{AbstractSimulationStepper} provides common evaluation methods and the pre/post listener hooks around them.

\section{Mass properties and rigid-body aggregation}
Component classes calculate local mass and centre of mass from geometry and materials. \class{MassCalculation} traverses the active configuration and its instances, translates local centres into assembly coordinates, respects applicable overrides, and aggregates the contributions. \class{MassCalculator} exposes convenient entry points for structure, launch, burnout, and time-dependent motor mass. \class{RigidBody} carries the resulting mass, centre, and inertias.

For ordinary contributions, the conceptual aggregation is
\begin{equation}
 m=\sum_i m_i,\qquad \boldsymbol r_{CG}=\frac{\sum_i m_i\boldsymbol r_i}{m},
 \qquad I=\sum_i\left(I_i+m_i d_i^2\right),
\end{equation}
with the appropriate axis and transformed instance positions. Overrides alter which contributions are included and how they are represented; they must be inspected before applying this expression to a particular design. The implementation distinguishes inertia about the longitudinal spin axis from transverse/pitch inertia. The flight stepper uses these effective inertias rather than a general deformable-body model.

Structure and motor contributions are evaluated separately. The structure normally remains constant within an unchanging active-stage configuration, while motor mass/CG follow motor time and stage activity. Staging changes the active set. A simulation actuator can also trigger component changes through the normal model API. Caches and change identifiers must therefore be consistent with both static design edits and supported live actuator properties.

\section{Motor data, installation, and ignition state}
\class{Motor} describes a motor model. \class{ThrustCurveMotor} supplies sampled thrust, mass/CG information, dimensions, manufacturer/designation metadata, and supported delays. It interpolates time-series values from the stored curve. The numerical integrator can evaluate a curve at intermediate Runge--Kutta stages; the motor is not advanced by incrementing a mutable sample index.

\class{MotorConfiguration} describes installation in a mount for a flight configuration, including ignition behavior and ejection delay. \class{MotorClusterState} is the per-run state for that installation: armed, thrusting, delaying, or spent, with actual ignition/cutoff/ejection times. Motor time is measured relative to ignition. Cluster count comes from the mount's physical instances and scales thrust appropriately.

The engine schedules ignition, burnout, and ejection events. An ejection charge is an event source that can activate ordinary recovery rules; it is not itself a deployed parachute. Plugged motors do not create the same ejection path. The MIT engine also schedules thrust-curve sample boundaries; FC-aware timing uses ignition-relative offsets so delayed ignition does not shift the curve into the wrong clock frame.

RASP \code{.eng} and RockSim motor loaders populate the motor database separately from rocket geometry. Unresolved references can generate loading warnings. Reproducibility requires curve identity/data and delay, not just a motor display label.

\section{Atmosphere and gravity}
\class{AtmosphericModel} returns \class{AtmosphericConditions} at altitude. The standard implementation, \class{ExtendedISAModel}, supports a layered ISA-derived atmosphere anchored either to standard conditions or custom launch pressure and temperature. It inherits interpolation support from \class{InterpolatingAtmosphericModel}. Atmospheric conditions supply density, sound speed, and viscosity-related quantities used by the aerodynamic path.

Pressure is in pascals and temperature in kelvin in the physical atmosphere. The FC adapter deliberately converts those values to hPa and Celsius. Launch altitude is mean-sea-level altitude; trajectory height above the pad is a different quantity. A pressure-derived barometric height zeroed during preflight is different again. These distinctions matter when comparing recorded telemetry with simulator truth.

\class{GravityModel} supplies gravity at a world location; \class{WGSGravityModel} is the normal implementation. \class{GeodeticComputationStrategy} converts local displacements to world coordinates and supplies the appropriate Coriolis term for the selected strategy. Flat, spherical, and other provided strategies are model choices, not display-only options.

\section{Wind and stochastic inputs}
\class{WindModel} is queried by time and altitude and can use MSL or AGL altitude references. \class{PinkNoiseWindModel} represents a baseline wind with correlated turbulence. \class{MultiLevelPinkNoiseWindModel} supports altitude-dependent profiles. Seeds are part of reproducibility: changing unrelated evaluation order should not casually introduce new random draws into a physical model.

The stepper constructs its relative-airflow vector using the wind model's convention. In the inspected code this calculation is \code{rocketVelocity.add(windVelocity)}, followed by rotation into the body frame. An external wind-vector source must be adapted to this convention rather than substituted under an assumed textbook sign. Bearing is clockwise from north; world horizontal axes are east and north.

Ensemble wind variation is layered onto these models: a per-flight bearing rotation and east/north offsets modify the existing profile while the baseline turbulence seed remains fixed across members. This is distinct from FC sensor noise, which perturbs the observations of the resulting flight.

\section{Coordinate and unit contracts}
\begin{tabularx}{\linewidth}{@{}L{3.3cm}YY@{}}\toprule
Boundary & Convention & Consequence\\\midrule
Component geometry & Longitudinal model X; placement relative to parent/assembly & CP and CG axial coordinates are model positions, not world altitude.\\
Flight world & X east, Y north, Z upward & Trajectory translation and world angular vectors use this frame.\\
Flight body & Longitudinal Z; quaternion maps body to world & Thrust, rotational rates, and sensor adapters require explicit transforms.\\
FC sensor mounting & Baseline X = body Z, Y = body X, Z = body Y & Longitudinal acceleration feeds sensor X; roll integration uses negative sensor X rate.\\
Display and persistence & Unit-group conversions; some file angles stored in degrees & Never copy displayed values directly into radian-valued setters.\\\bottomrule
\end{tabularx}

\begin{contract}
These are rigid-body engineering models, without fluid-structure, combustion, flexible-body, or complete electronics dynamics. A replacement submodel must preserve units, frames, timing, and listener contracts.
\end{contract}
\sourceline{\source{core/src/main/java/info/openrocket/core/masscalc/MassCalculation.java}{MassCalculation}; \source{core/src/main/java/info/openrocket/core/masscalc/RigidBody.java}{RigidBody}; \source{core/src/main/java/info/openrocket/core/motor/ThrustCurveMotor.java}{ThrustCurveMotor}; \source{core/src/main/java/info/openrocket/core/simulation/MotorClusterState.java}{MotorClusterState}; \source{core/src/main/java/info/openrocket/core/models/atmosphere/ExtendedISAModel.java}{ExtendedISAModel}; \source{core/src/main/java/info/openrocket/core/util/GeodeticComputationStrategy.java}{GeodeticComputationStrategy}.}

\section{How shape becomes mass, CG, and inertia}
For a constant-radius hollow tube, the structural volume is conceptually $\pi(r_o^2-r_i^2)\ell$, multiplied by bulk material density. Its CG lies halfway along its uniform length, and hollow-cylinder inertia formulas distinguish rotation about the long axis from rotation about a transverse axis. Internal rings use the same kind of volume accounting with their own dimensions. A mass object instead supplies explicit mass and an approximate envelope for inertia.

Curved \class{SymmetricComponent} shapes use 128 axial divisions. At each division the code evaluates the external radius at both ends and treats the segment as a conical frustum. A hollow shell subtracts an inner frustum, using wall slope to convert specified thickness into radial material depth; a filled component uses zero inner radii. Segment volumes, first moments, and axial/transverse inertias are accumulated, followed by the parallel-axis shift to the component CG.

The same traversal separately accumulates full external volume, wetted surface area, side-projected planform area, and planform centroid. Full volume describes the outer shape regardless of whether the part is hollow. It is used in aerodynamic CP relationships. Material volume describes what is actually made of the selected material. Substituting one for the other would make a hollow nose's aerodynamic CP incorrectly depend on its wall thickness.

Fin properties are reduced from the outline and supporting root. Material volume includes the profile's thickness factor; mounting tabs and fillets add their own volumes and centroids. Repeated fins, rings of pods, and motor clusters require transformed instance locations before aggregation. Moving a component sideways or rotating an assembly can affect inertia even if total mass and axial CG barely change.

Mass and CG overrides operate at their selected scope. A component-only mass override retains its children as separate contributions. A subtree mass override suppresses the children's aggregate mass contribution; a subtree CG override relocates the children's aggregate center according to the code's override semantics. Unit-inertia approximations remain tied to component geometry. These options are useful for reconciling a measured vehicle, but they are not a procedure for identifying a full measured inertia tensor.

\sourceline{\source{core/src/main/java/info/openrocket/core/rocketcomponent/SymmetricComponent.java}{calculateProperties}; \source{core/src/main/java/info/openrocket/core/rocketcomponent/FinSet.java}{fin mass, tabs, and fillets}; \source{core/src/main/java/info/openrocket/core/masscalc/MassCalculation.java}{structure/motor traversal and overrides}; \source{core/src/main/java/info/openrocket/core/masscalc/MassCalculator.java}{mass calculation entry points}.}

\section{Motor curves, propulsion timing, and database identity}
A thrust-curve motor stores paired time/thrust samples and mass/CG information. \class{ThrustCurveMotor} locates the surrounding sample interval and linearly interpolates thrust and center-of-mass data at the requested motor time. The same object can be queried at several provisional solver times without changing its state. Calculated metadata such as total impulse, average thrust, maximum thrust, and estimated burn duration supports selection and reporting; the designation is not a replacement for the curve.

Loaders handle the different information available in RASP and RockSim motor formats. They validate and normalize input, and may derive mass evolution when the source lacks a full measured mass history. Manufacturer normalization, motor digests, curve sets, and the database finder help distinguish similar names and alternate measurements. A nominal ``same motor'' selection can still refer to different curve data, so reproduction should retain the actual resolved identity.

The installation's \class{IgnitionEvent} can be launch, the lower stage's ejection or burnout, automatic, or never, with an additional configured delay. Automatic selects launch for a launch stage and ejection-triggered ignition for an upper stage. The running \class{MotorClusterState} tracks when that condition actually occurred. A motor's thrust and remaining mass use elapsed time since its own ignition, while the flight branch uses simulation time. Delay/ejection and stage separation are distinct event paths.

Thrust is summed over active motor installations and their physical instances. The model reads supplied thrust curves; it does not solve chamber pressure, nozzle flow, combustion instability, or an electrical ignition circuit. A cluster configuration repeats motor geometry and thrust; it does not imply a modeled ignition-failure distribution unless an explicit extension supplies one.

\sourceline{\source{core/src/main/java/info/openrocket/core/motor/ThrustCurveMotor.java}{ThrustCurveMotor}; \source{core/src/main/java/info/openrocket/core/motor/IgnitionEvent.java}{IgnitionEvent}; \source{core/src/main/java/info/openrocket/core/database/motor/ThrustCurveMotorSet.java}{ThrustCurveMotorSet}; \source{core/src/main/java/info/openrocket/core/simulation/MotorClusterState.java}{MotorClusterState}.}

\section{Atmospheric quantities used by the solver}
The atmosphere supplies pressure and temperature; \class{AtmosphericConditions} derives dry-air density as $\rho=P/(RT)$ with $R=287.053\ \mathrm{J\,kg^{-1}K^{-1}}$. In this implementation the sound speed is the linear approximation $a=165.77+0.606T$ in m/s, with $T$ in kelvin. Kinematic viscosity is $(3.7291\times10^{-6}+4.9944\times10^{-8}T)/\rho$. It is therefore inaccurate to describe this code as evaluating the exact ideal-gas sound-speed expression or full Sutherland viscosity at every sample.

\class{ExtendedISAModel} uses layer boundaries at 0, 11, 20, 32, 47, 51, 71, and 84.852 km, with prescribed base temperatures and hydrostatic pressure relations. Custom launch conditions anchor the lower atmosphere and connect to the higher layer structure. Its constructor restricts a custom first altitude to below 11 km and validates positive absolute temperature and pressure. An interpolation wrapper caches sampled atmospheric values for repeated queries.

The standalone \class{AtmosphericConditions} default is 293.15 K, whereas the standard \class{ExtendedISAModel} sea-level baseline is 288.15 K. Code that constructs a fresh conditions object for a static analysis is not necessarily using the same atmosphere as a flight simulation at its configured launch site. The source also records known approximation limits at high altitude. Humidity, precipitation, and detailed meteorological assimilation are not modeled by these classes.

\section{Wind profiles, gravity, and Earth geometry}
The average wind model combines mean speed/direction with seeded correlated pink-noise turbulence. Turbulence intensity and standard deviation describe variability around the mean; they are not a sequence of independent draws at each RK evaluation. The multilevel model stores altitude-dependent wind levels and interpolates profile behavior, with an explicit AGL/MSL altitude reference. Wind-profile import and editing prepare this model before execution.

Gravity uses a latitude-dependent WGS-style surface magnitude and an inverse-square altitude correction. The geodetic strategy separately controls how local displacement becomes geographic position: \code{FLAT} uses meters-per-degree scaling and zero Coriolis acceleration; \code{SPHERICAL} uses great-circle geometry; \code{WGS84} uses an ellipsoid/Vincenty direct solution. The latter two include the implementation's Coriolis term. Selecting a more detailed map geometry does not replace the rest of the force model.

\sourceline{\source{core/src/main/java/info/openrocket/core/models/atmosphere/AtmosphericConditions.java}{AtmosphericConditions}; \source{core/src/main/java/info/openrocket/core/models/atmosphere/InterpolatingAtmosphericModel.java}{InterpolatingAtmosphericModel}; \source{core/src/main/java/info/openrocket/core/models/wind/PinkNoiseWindModel.java}{PinkNoiseWindModel}; \source{core/src/main/java/info/openrocket/core/models/wind/MultiLevelPinkNoiseWindModel.java}{MultiLevelPinkNoiseWindModel}; \source{core/src/main/java/info/openrocket/core/models/gravity/WGSGravityModel.java}{WGSGravityModel}.}

